\section{Threat Hunting Fundamentals}

Threat hunting is the proactive search for adversary activity that has evaded automated detection. Unlike alert triage, hunting begins without a predefined signal; instead, the analyst forms hypotheses based on telemetry, threat intelligence, and knowledge of system behavior. The goal is to identify subtle patterns, weak signals, and behavioral anomalies that indicate malicious activity. This section introduces the core mindset and techniques used in modern threat hunting, supported by SPL and KQL examples that reflect real investigative workflows.

A hunt typically begins with a hypothesis. This can be driven by CTI (“a threat actor is known to use encoded PowerShell commands”), by anomalies in telemetry (“a spike in failed logons”), or by environmental knowledge (“this server should never initiate outbound connections”). Once a hypothesis is defined, the hunter uses SPL or KQL to explore relevant telemetry, pivot across related events, and refine the search until meaningful patterns emerge. This iterative process is what separates hunting from simple querying.

One of the most common starting points in a hunt is process creation activity. Sysmon Event ID 1 provides detailed insight into how processes are launched, including parent-child relationships and command-line arguments. The following example searches for suspicious PowerShell usage, a frequent indicator of malicious activity:

\textbf{SPL Example}
\begin{lstlisting}
index=sysmon EventCode=1 Image="*powershell.exe"
| search CommandLine="*EncodedCommand*"
| table _time, Computer, User, ParentImage, Image, CommandLine
\end{lstlisting}

\textbf{KQL Equivalent}
\begin{lstlisting}
Sysmon
| where EventID == 1 and Image endswith "powershell.exe"
| where CommandLine contains "EncodedCommand"
| project TimeGenerated, Computer, User, ParentImage, Image, CommandLine
\end{lstlisting}

Another common hunting technique is identifying rare or unusual behavior. Adversaries often stand out not because their activity is obviously malicious, but because it is uncommon compared to normal system patterns. The following example highlights rare parent-child process pairs:

\textbf{SPL Example}
\begin{lstlisting}
index=sysmon EventCode=1
| stats count by ParentImage, Image
| where count < 3
\end{lstlisting}

\textbf{KQL Equivalent}
\begin{lstlisting}
Sysmon
| where EventID == 1
| summarize Count=count() by ParentImage, Image
| where Count < 3
\end{lstlisting}

Authentication activity is another rich source of hunting leads. Repeated failed logons, unusual logon types, or logons from unexpected hosts can indicate credential misuse or lateral movement. The following example identifies accounts with excessive failed logons:

\textbf{SPL Example}
\begin{lstlisting}
index=wineventlog EventCode=4625
| stats count by AccountName, WorkstationName
| where count > 10
\end{lstlisting}

\textbf{KQL Equivalent}
\begin{lstlisting}
SecurityEvent
| where EventID == 4625
| summarize Count=count() by Account, Workstation
| where Count > 10
\end{lstlisting}

Network telemetry also plays a critical role in hunting. DNS logs, proxy logs, and NetFlow data can reveal command-and-control activity, data exfiltration, or lateral movement. The following example identifies rare outbound connections:

\textbf{SPL Example}
\begin{lstlisting}
index=sysmon EventCode=3
| stats count by DestinationIp, DestinationPort
| where count < 5
\end{lstlisting}

\textbf{KQL Equivalent}
\begin{lstlisting}
Sysmon
| where EventID == 3
| summarize Count=count() by DestinationIp, DestinationPort
| where Count < 5
\end{lstlisting}

Threat hunting is ultimately about pattern recognition and curiosity. The examples in this section demonstrate the foundational techniques—filtering, aggregation, rarity analysis, and behavioral correlation—that underpin more advanced hunts. As the handbook progresses, these fundamentals will be expanded into full hunting playbooks and detection patterns that reflect real-world adversary behavior.
